%%%% fatec-article.tex, 2024/03/10

%% Classe de documento
\documentclass[
  a4paper,%% Tamanho de papel: a4paper, letterpaper (^), etc.
  12pt,%% Tamanho de fonte: 10pt (^), 11pt, 12pt, etc.
  english,%% Idioma secundário (penúltimo) (>)
  brazilian,%% Idioma primário (último) (>)
]{article}

%% Pacotes utilizados
\usepackage[]{fatec-article}

\Author{1}{Name={Lucas de Barros Parada\\Paulo Seiji Yamamoto Junior\\Filipe de Oliveira Lima\\Bruno de Lucca Soares Miranda}}

\Author{2}{Name={\{ lucas.parada@fatec.sp.gov.br \}\\ \{ paulo.yamamoto01@fatec.sp.gov.br \} \\ \{ filipe.lima18@fatec.sp.gov.br\} \\ \{ bruno.miranda19@fatec.sp.gov.br \}}}

%% Definição das palavras-chaves/keywords
\Keyword{1}{Aprendizagem Profunda}{Deep Learning}
\Keyword{2}{Cogumelos}{Mushrooms}
\Keyword{3}{Internet das Coisas}{Internet of Things}

%%%% Resumo no idioma primário (brazilian)
\begin{Abstract}[brazilian]%% Idioma (brazilian ou english)
O aumento no número de infecções em plantações de cogumelos tem levado produtores a investir em estratégias mais eficazes de monitoramento e controle. Diante desse cenário, este projeto propõe o desenvolvimento de um sistema inteligente voltado ao acompanhamento do cultivo de cogumelos, integrando sensores ambientais e técnicas de inteligência artificial. O sistema realiza o monitoramento em tempo real de variáveis como umidade, temperatura e possíveis focos de infecção, utilizando câmeras para detectar anomalias visuais, como manchas e alterações de cor. Para aprimorar a precisão das análises, o sistema emprega algoritmos de aprendizagem profunda, capazes de reconhecer padrões complexos e gerar resultados mais confiáveis sobre o estado das culturas. Ao identificar sinais de contaminação, o sistema emite alertas automáticos e fornece suporte à tomada de decisão do produtor. Dessa forma, busca-se reduzir perdas, otimizar o manejo produtivo e fomentar a inovação tecnológica no setor da fungicultura.
\end{Abstract}

%%%% Resumo no idioma secundário (english)
\begin{Abstract}[english]%% Idioma (brazilian ou english)
The increase in infections in mushroom crops has led producers to invest in more effective monitoring and control strategies. In this context, this project proposes the development of an intelligent system aimed at monitoring mushroom cultivation by integrating environmental sensors and artificial intelligence techniques. The system performs real-time monitoring of variables such as humidity, temperature, and possible infection sources, using cameras to detect visual anomalies such as spots and color changes. To improve analytical accuracy, the system employs deep learning algorithms capable of recognizing complex patterns and generating more reliable results regarding the condition of the crops. Upon detecting signs of contamination, the system issues automatic alerts and supports the producer in decision-making. Thus, it seeks to reduce losses, optimize production management, and promote technological innovation in the mushroom cultivation sector.
\end{Abstract}

%% Processamento de entradas (itens) do índice remissivo (makeindex)
\makeindex%

%% Arquivo(s) de referências
\addbibresource{Ref.bib}

%% Início do documento
\begin{document}

% Seções e subseções

%\section{}%

%\subsection{Título de Seção Secundária}%

%\subsubsection{Título de Seção Terciária}%

%\paragraph{Título de seção quaternária}%

%\subparagraph{Título de seção quinária}%

\section*{Introdução}%
\label{sect:intro}
Os cogumelos constituem uma relevante fonte natural de nutrientes, destacando-se pelo elevado teor de proteínas, vitaminas e compostos bioativos benéficos à saúde humana. No mercado global de cogumelos comestíveis, as doenças fúngicas representam uma das principais causas de perdas produtivas. Esses organismos são cultivados predominantemente em ambientes fechados, sob condições ambientais controladas, o que favorece a adoção de programas integrados de manejo de doenças. Tais programas combinam o uso de fungicidas químicos, alternativas de biocontrole e práticas agronômicas adequadas, permitindo reduzir a incidência de patógenos e garantir o manejo sanitário das culturas. Dessa forma, diversos métodos de controle são continuamente aplicados com o objetivo de minimizar perdas e assegurar a qualidade do produto final.\cite{Bellettini2017}

Com base nesse cenário, a aplicação de tecnologias emergentes no setor agrícola possibilita o aumento da eficiência produtiva, a redução de perdas e a melhoria da qualidade dos alimentos e fatores diretamente relacionados à segurança alimentar. Além disso, o monitoramento contínuo de variáveis ambientais, como umidade e temperatura, permite otimizar o uso de recursos naturais, reduzir o desperdício de insumos e promover uma produção mais sustentável, resiliente e alinhada às demandas atuais por sistemas agrícolas inteligentes e ecologicamente responsáveis.

No contexto brasileiro, o Pleurotus ostreatus (shimeji) destaca-se entre as espécies mais amplamente cultivadas, devido ao seu elevado valor nutricional e à crescente demanda no mercado alimentício. Entretanto, o aumento da produção tem sido acompanhado por uma maior ocorrência de contaminações microbiológicas, resultantes da proliferação de microrganismos oportunistas durante as etapas de cultivo e processamento, podendo ocasionar perdas totais nas colheitas. As infestações por pragas e doenças comprometem significativamente a qualidade e a composição do produto final, reduzindo a competitividade do produtor no mercado e impactando negativamente a rentabilidade e a sustentabilidade da produção.\cite{Scheffer2023}

Com isso o Vale do Ribeira, localizado no Estado de São Paulo, destaca-se como uma região promissora para a implementação de sistemas baseados em tecnologias inteligentes. Essa região, foco da presente pesquisa, apresenta um constante crescimento na produção de espécies de cogumelos, especialmente o shimeji, o que reforça o potencial de aplicação de soluções tecnológicas voltadas ao monitoramento e controle de variáveis ambientais, contribuindo para o fortalecimento da produção local e o desenvolvimento sustentável do setor.\cite{Urben2017}

Deste modo o presente projeto tem como objetivo o desenvolvimento de um sistema inteligente de monitoramento do cultivo de cogumelos, voltado à prevenção de infecções e infestações por parasitas. A proposta visa integrar sensoriamento ambiental e inteligência artificial (IA) para garantir o acompanhamento automatizado das condições do composto e do desenvolvimento dos cogumelos. O sistema será responsável por coletar e exibir dados em tempo real sobre umidade, temperatura e porcentagem estimada de infecção, obtidos por meio de sensores instalados no ambiente de cultivo, permitindo que o produtor acompanhe remotamente as condições ideais para o crescimento saudável dos fungos.

Além disso, será implementada uma IA baseada em aprendizagem profunda (deep learning). As imagens capturadas no ambiente serão processadas e utilizadas tanto para o treinamento da IA quanto para o monitoramento contínuo do cultivo. O modelo será treinado para identificar padrões visuais indicativos de infecções e parasitas, como alterações de cor, presença de manchas e irregularidades no composto. Quando o sistema detectar anomalias que indiquem risco de contaminação, a porcentagem de alerta será automaticamente ajustada e o produtor será notificado por meio do aplicativo, permitindo uma resposta rápida e precisa no manejo das culturas. 

Alinhando-se o presente estudo com a Agenda 2030 e os Objetivos de Desenvolvimento Sustentável (ODS) propostos pela Organização das Nações Unidas (ONU), a presente pesquisa relaciona-se diretamente ao ODS nº 2 — Fome Zero e Agricultura Sustentável. Este objetivo visa erradicar a fome e promover sistemas alimentares produtivos, resilientes e ambientalmente responsáveis. Dessa forma, o desenvolvimento de um sistema inteligente de monitoramento e prevenção no cultivo de cogumelos contribui não apenas para o aumento da produtividade e da segurança alimentar, mas também para a redução do desperdício e o uso eficiente dos recursos naturais, reforçando o compromisso com práticas agrícolas sustentáveis e inovadoras.


\section*{OBJETIVOS} \label{sect:obj}

Com esta pesquisa, tem-se como principal objetivo Desenvolver um sistema inteligente de monitoramento para o cultivo de cogumelos, utilizando sensores ambientais e inteligência artificial, a fim de prevenir infecções e infestações por parasitas, proporcionando maior controle, segurança e produtividade ao produtor, utilizando níveis de umidade, temperatura e índice de infecção. Além disso, como metas parciais visa-se:

- Realizar o treinamento do modelo de aprendizagem profunda (Deep Learning) visando alcançar uma acurácia mínima de 50\% e uma taxa de falsos positivos inferior a 10\%, assegurando desempenho satisfatório na detecção e classificação dos dados analisados. 

- Implementar uma integração eficiente entre o sistema proposto e a área de atuação dos produtores do Vale do Ribeira, ampliando o acesso e a disponibilidade de dados para o aprimoramento contínuo do modelo de inteligência artificial. 

- Desenvolver a integração de câmeras e sensores de temperatura e umidade em um ambiente de Internet das Coisas (IoT) funcional, possibilitando a coleta automatizada de dados e o monitoramento em tempo real das condições ambientais do cultivo. 

\section*{ESTADO DA ARTE} \label{sect:estadoarte}

Nos últimos anos, o uso de técnicas de machine learning (ML) e visão computacional na agricultura tem se expandido significativamente, oferecendo soluções avançadas para monitoramento, classificação e colheita de produtos cultivados em ambientes controlados. Em especial para este estudo, a aplicação dessas técnicas no cultivo de cogumelos tem se mostrado promissora. 

Diversos estudos têm explorado algoritmos de detecção de objetos para localizar cogumelos em imagens e determinar o estágio de crescimento, visando otimizar a colheita.\cite{MOYSIADIS2023} utilizaram o algoritmo YOLOv5 para classificar cogumelos em três estágios de crescimento, chegando a obter uma acurácia de 70\% para o estágio final, quando os cogumelos estão prontos para a colheita. O estudo demonstrou que a classificação do sistema pode contribuir para decisões mais eficientes sobre o momento ideal, evitando perdas por amadurecimento excessivo ou coleta precoce. 

Outros estudos relevantes incluem a utilização do SSD (Single Shot MultiBox Detector) em robôs colhedores de cogumelos, que alcançaram até 95\% de acurácia na detecção e 86,8\% de sucesso na colheita individual, evidenciando a viabilidade prática de sistemas automatizados em ambientes controlados. Pesquisas voltadas para a classificação de espécies de cogumelos comestíveis e venenosos mostraram que ensembles de classificadores (Decision Tree, KNN, Random Forest, SVM, Naive Bayes) podem atingir 100\% de acurácia, reforçando a importância da inteligência artificial na diferenciação de cogumelos comestíveis e tóxicos. 

A integração de Internet das Coisas (IoT) com ML em sistemas de fungicultura inteligente também tem recebido destaque. Sensores de temperatura e umidade, associados a microcontroladores, permitem ajustar automaticamente variáveis críticas para o crescimento do micélio e do corpo de frutificação. Algoritmos de ML garantem a identificação precisa de cogumelos comestíveis, prevenindo a produção e consumo de exemplares tóxicos. Resultados experimentais indicam que esses sistemas mantêm condições ideais de temperatura e umidade, reduzem o esforço manual e fornecem dados em tempo real para os produtores. Avaliações de usabilidade (SUS) também demonstraram alto nível de satisfação, com pontuação média de 82\%, evidenciando a aplicabilidade prática dessas soluções. 

De forma complementar, \cite{Rahman2021} demonstrou que a integração de sensores ambientais com sistemas IoT permite monitoramento contínuo em tempo real e emissão de alertas automáticos, garantindo respostas rápidas a alterações ambientais e promovendo uma gestão mais eficiente do processo produtivo.  

Além do monitoramento, a classificação de espécies de cogumelos comestíveis e venenosos é um foco relevante na literatura. \cite{MetlekCetiner2023} desenvolveram um estudo utilizando um conjunto de dados com 8.124 amostras e 22 atributos visuais e físicos. Foram aplicados os algoritmos Random Forest (RF), Decision Tree (DT) e Regressão Logística (LR), com otimização de hiperparâmetros via GridSearchCV. Os resultados demonstraram que o RF apresentou melhor desempenho, alcançando precisão, recall e F1-score de até 0,97, 0,98 e 0,95, respectivamente, enquanto a DT otimizada destacou-se na distinção da classe venenosa. A regressão logística apresentou desempenho competitivo, mas ligeiramente inferior. O estudo evidencia que a seleção adequada de atributos e a otimização de parâmetros são fundamentais para reduzir erros na identificação de espécies. 

Dessa forma, a integração entre algoritmos de aprendizagem profunda para análise visual, sensoriamento ambiental e tecnologias de Internet das Coisas (IoT) para monitoramento em tempo real configura uma abordagem tecnológica robusta e de elevada relevância para a produção de cogumelos. Soluções convencionais frequentemente apresentam restrições de acesso para usuários e produtores em campo, decorrentes de limitações de infraestrutura e da complexidade operacional. O sistema proposto neste projeto trata tais barreiras, viabilizando monitoramento remoto contínuo e fornecendo suporte automatizado à tomada de decisão. Essa abordagem estabelece uma base sólida para o desenvolvimento de sistemas inteligentes aplicáveis diretamente ao cultivo de espécies como o shimeji, promovendo incremento na produtividade, melhoria na segurança alimentar e maior sustentabilidade da fungicultura.

\section*{METODOLOGIA} \label{sect:metodologia}

Para a construção do sistema proposto, será implementada uma arquitetura de inteligência artificial baseada em aprendizado profundo (Deep Learning), atuando como núcleo analítico do processo de monitoramento e detecção automatizada. O sistema integrará sensores ópticos (câmeras com microcontrolador integrado) e sensores ambientais de temperatura, umidade relativa e luminosidade, interligados por uma infraestrutura IoT (Internet of Things), responsável pela aquisição contínua e transmissão dos dados provenientes das estufas de cultivo.

Durante o processo de coleta, múltiplos fluxos de imagem (bloco 1) e dados ambientais (bloco 2) serão capturados em tempo real, encaminhados para o sistema via Wi-Fi (blocos 3 e 4), indexados e armazenados em um banco de dados não relacional (blocos 5 e 6), garantindo integridade e disponibilidade para uso posterior em tarefas de treinamento e validação de modelos.

O conjunto de imagens coletadas comporá o dataset primário para o treinamento supervisionado da IA. O modelo de aprendizado profundo, baseado em redes neurais, será configurado para detecção de padrões morfológicos e análises de variações de coloração e textura, possibilitando a classificação automatizada dos estágios de crescimento e a identificação precoce de anomalias nos compostos.

Essa abordagem possibilita a criação de um sistema autônomo e adaptativo, capaz de realizar inferência em tempo real, monitorar o desenvolvimento dos cogumelos com alta precisão e fornecer alertas preditivos ao produtor sobre possíveis contaminações ou irregularidades no ambiente de cultivo.

\begin{figure}[!ht]
\centering
\SetCaptionWidth{\ifbool{@LayoutA}{0.7}{0.72}\linewidth}
\caption{Fluxograma da metodologia adotada para desenvolvimento do sistema para detecção de parasitas em cogumelos}
\includegraphics[width = 450px]{fluxograma.drawio.png}
\SourceOrNote{Autoria própria, 2026}
\end{figure}

Durante o funcionamento em ambiente real, o sistema executa o processamento em tempo real de imagem, capturando aproximadamente 1 frame por minuto (bloco 7). Cada frame é enviado ao modelo de IA, que realiza a inferência instantânea e gera uma estimativa probabilística de contaminação com base em características morfológicas detectadas. Quando um evento anômalo é identificado — como a presença de fungos invasores ou alterações visuais atípicas — o sistema gera alertas automáticos (bloco 8), notificando o produtor por meio do aplicativo integrado. Da mesma forma, os dados de avaliação da imagem são salvos no banco de dados. O software foi projetado para operar em uma aplicação mobile, destinada ao monitoramento em tempo real do ambiente de cultivo, exibição de alertas, acompanhimento de métricas de umidade, temperatura e iluminação, e consulta de painéis analíticos dinâmicos com gráficos de desempenho.

A comunicação entre as instâncias é realizada via API RESTful, permitindo sincronização contínua de dados e emissão de relatórios automáticos com informações sobre a saúde dos compostos, eficiência ambiental e status de contaminação. Essa integração garante que o produtor disponha de feedbacks em tempo real sobre o estado de suas produções e insumos preditivos para otimização de recursos e prevenção de perdas.

Como resultado, o sistema fornece uma plataforma inteligente e autônoma de monitoramento preditivo de fungicultura, combinando visão computacional, IoT e aprendizado profundo em uma arquitetura escalável, voltada à melhoria da produtividade, qualidade e sustentabilidade dos cultivos de Pleurotus ostreatus (shimeji).













\printbibliography


%% Elementos pós-textuais (opcionais): Apêndice e Anexo
%Caso for utilizar, basta retirar o símbolo de % na frente do comando
%\input{./Extras/post-textual}

%% Fim do documento
\end{document}